\chapter{Phụ Lục}
\label{chap:appendix}

% ─────────────────────────────────────────────────────────────────────────────
\section{Cấu Trúc Mã Nguồn}
\label{apx:repo}

\begin{lstlisting}[language=, numbers=none]
it-interview-recsys/
|-- app.py                  # Streamlit shell + FastAPI bridge + sinh static
|-- config.py               # LLM model, KC_BY_ROLE (17 KC), API_PORT
|-- data/
|   |-- raw/question_bank/  # question_bank.json (2122 cau)
|   +-- schemas/            # JSON Schema: question, taxonomy, JD, profile
|-- src/
|   |-- data/   # data_loader, question_manager, schema_validator
|   |-- kt/     # competency_engine, recommender, recommendation_engine,
|   |           # scoring, kt_predictor, kt_models
|   |-- ai/     # grader, prompts
|   +-- api/    # serving (FastAPI), user_manager, test_generator
|-- app/frontend/           # JSX (no bundler): state, chrome, widgets,
|                           # screens-*, router + styles.css
|-- kt_models/              # BKT, DKT, SAKT (kien truc mo hinh)
|-- notebooks/              # 01_eda -> 02_feature -> 03_train -> 04_eval
|-- results/                # models/, plots/, metrics/, ablation/
+-- utils/data/             # collection/ generation/ qc/ processing/ simulation/
\end{lstlisting}

% ─────────────────────────────────────────────────────────────────────────────
\section{Prompt Chấm Điểm (Gemini)}
\label{apx:prompt}

Prompt dùng để chấm câu tự luận/coding bằng Gemini (\texttt{src/ai/prompts.py}). Các biến trong dấu ngoặc nhọn được điền tại thời điểm chạy:

\begin{center}
\fbox{\parbox{0.95\linewidth}{\small
Bạn là chuyên gia chấm điểm phỏng vấn kỹ thuật. Hãy chấm điểm khách quan và nhất quán.\\[3pt]
\textbf{Câu hỏi:} \{question\}\\
\textbf{Đáp án tham khảo:} \{expected\_answer\}\\
\textbf{Các ý chính cần có:} \{evaluation\_points\}\\
\textbf{Câu trả lời ứng viên:} \{candidate\_answer\}\\
\textbf{Rubric:} \{scoring\_rubric\}\\[3pt]
\textbf{Thang điểm:}\\
0--3: Sai hoặc không liên quan\\
4--5: Đúng một phần, thiếu ý chính\\
6--7: Đủ ý chính, thiếu chiều sâu / ví dụ\\
8--9: Đầy đủ, có ví dụ cụ thể, giải thích rõ\\
10: Xuất sắc --- đề cập trade-off và edge case\\[3pt]
Trả về JSON hợp lệ: \texttt{\{"score": <0-10>, "feedback": "...", "strengths": [...], "improvements": [...]\}}
}}
\end{center}

Khi gọi LLM thất bại, hệ thống lui về \textbf{local rubric}: một \texttt{evaluation\_point} được coi là ``đạt'' khi $\geq 50\%$ từ khóa của nó xuất hiện trong câu trả lời.

% ─────────────────────────────────────────────────────────────────────────────
\section{Đặc Tả API (chi tiết)}
\label{apx:api}

Ví dụ \texttt{POST /api/grade}:
\begin{lstlisting}[language=, numbers=none]
// Request
{ "question": { "question_id": "...", "question_type": "THEORY", ... },
  "candidate_response": { "answer": "...", "code": "" },
  "user_id": "u_123" }

// Response
{ "score": 0.0-1.0, "is_correct": true|false|null,
  "feedback": "...", "strengths": [...], "improvements": [...],
  "method": "gemini" | "local_rubric",
  "gemini_error": "<if any>" }
\end{lstlisting}

Ví dụ \texttt{POST /api/assessment/finalize}: nhận \texttt{stage1\_results}, \texttt{stage2\_results}, chạy KT pipeline (EMA + blend) và trả về skill vector cập nhật cùng phân loại trình độ.

% ─────────────────────────────────────────────────────────────────────────────
\section{Siêu Tham Số Mô Phỏng và Mô Hình}
\label{apx:hparams}

\begin{table}[H]
  \centering
  \caption{Tham số pipeline mô phỏng dữ liệu}
  \begin{tabular}{ll}
    \toprule
    \textbf{Tham số} & \textbf{Giá trị} \\
    \midrule
    Số người dùng ($N$)            & 1.000 (DA 334 / DS 333 / DE 333) \\
    Số câu/người (thiết kế)        & 20--60 (thực tế cap $\approx$40) \\
    Noise IRT ($\sigma$)           & 0.30 \\
    Learning rate ($\eta$)         & 0.08 \\
    Forgetting rate ($\gamma$)     & 0.002 \\
    Difficulty $\to$ giá trị       & EASY 0.30 / MEDIUM 0.55 / HARD 0.80 \\
    Seed                           & 42 \\
    \bottomrule
  \end{tabular}
\end{table}

\begin{table}[H]
  \centering
  \caption{Hằng số runtime của hệ thống gợi ý}
  \begin{tabular}{ll}
    \toprule
    \textbf{Tham số} & \textbf{Giá trị} \\
    \midrule
    Blend KT / EMA ($\lambda$)        & 0.70 / 0.30 \\
    EMA smoothing ($\alpha$)          & 0.35 \\
    Lịch sử tối thiểu kích hoạt KT    & 3 tương tác \\
    Ngưỡng KC yếu / mạnh              & 0.40 / 0.70 \\
    Trọng số gợi ý (focus weak)       & 0.65 weak + 0.30 diff + 0.05 cov \\
    Khoảng dự đoán KT                 & [0.05, 0.99] \\
    \bottomrule
  \end{tabular}
\end{table}

% ─────────────────────────────────────────────────────────────────────────────
\section{Quy Trình Notebooks và Tái Lập}
\label{apx:reproducibility}

Phần huấn luyện Knowledge Tracing được tổ chức thành bốn notebook trong thư mục \texttt{notebooks/}, chạy \textbf{tuần tự}, mỗi notebook nhận đầu ra của notebook trước (tổng quan ở Mục~\ref{sec:notebook_pipeline}). Bảng~\ref{tab:notebook_io} tóm tắt đầu vào/đầu ra của từng notebook.

\begin{table}[H]
  \centering
  \caption{Đầu vào và đầu ra của bốn notebook}
  \label{tab:notebook_io}
  \begin{tabular}{p{2.5cm}p{4.2cm}p{4.6cm}}
    \toprule
    \textbf{Notebook} & \textbf{Đầu vào} & \textbf{Đầu ra} \\
    \midrule
    \texttt{01\_eda} &
    \texttt{data/raw/*} (question bank, virtual users, interaction log, self rating log) &
    Biểu đồ EDA trong \texttt{results/plots/} + danh sách issue cần fix \\
    \addlinespace
    \texttt{02\_feature\_engineering} &
    \texttt{interaction\_log.csv} (đã fix), \texttt{self\_rating\_log.csv} &
    \texttt{data/processed/} \{train,val,test\}\texttt{\_seqs.pkl} + encoder/metadata \\
    \addlinespace
    \texttt{03\_train\_models} &
    \texttt{data/processed/} \texttt{*\_seqs.pkl} &
    \texttt{results/models/} (\texttt{bkt.npy}, \texttt{dkt\_*.pt}, \texttt{sakt\_*.pt}), \texttt{results/metrics/}, \texttt{results/ablation/} \\
    \addlinespace
    \texttt{04\_evaluation} &
    \texttt{results/models/*} + \texttt{data/processed/*\_seqs.pkl} &
    Biểu đồ phân tích sâu trong \texttt{results/plots/}; kết luận chọn mô hình \\
    \bottomrule
  \end{tabular}
\end{table}

\noindent\textbf{Các bước chạy lại:}
\begin{enumerate}
  \item Đặt dữ liệu thô vào \texttt{data/raw/}.
  \item Chạy tuần tự: \texttt{01\_eda} $\to$ \texttt{02\_feature\_engineering} $\to$ \texttt{03\_train\_models} $\to$ \texttt{04\_evaluation}.
  \item Giữ \texttt{seed = 42} ở mọi notebook để tái lập kết quả; hai notebook \texttt{03}--\texttt{04} cần GPU (Kaggle Tesla T4, PyTorch 2.6).
  \item Mô hình tốt nhất (\texttt{dkt\_quality.pt}) được nạp ở runtime bởi \texttt{src/kt/kt\_predictor.py} cho hệ thống web (xem Chương~\ref{chap:system}).
\end{enumerate}

\noindent Tám cấu hình ablation trong \texttt{03\_train\_models} được liệt kê đầy đủ ở Bảng~\ref{tab:configs} (Chương~\ref{chap:experiment}). Trước khi huấn luyện, notebook kiểm tra toàn vẹn dữ liệu theo bảy tiêu chí: không rò rỉ user giữa train/test; phân phối nhãn train$\approx$test; độ dài chuỗi; drift tần suất KC; mất cân bằng lớp (dùng PR-AUC); phân phối quality; và độ phủ cold-start.

% ─────────────────────────────────────────────────────────────────────────────
\section{Năm Loại Câu Hỏi}
\label{apx:qtypes}

\begin{table}[H]
  \centering
  \caption{Các loại câu hỏi trong question bank}
  \begin{tabular}{lll}
    \toprule
    \textbf{Loại} & \textbf{Mô tả} & \textbf{Cách chấm} \\
    \midrule
    MC\_SINGLE  & Trắc nghiệm 1 đáp án & Tất định (so khớp option) \\
    TRUE\_FALSE & Đúng/Sai            & Tất định \\
    FILL\_BLANK & Điền chỗ trống      & Tất định (accepted answers) \\
    THEORY      & Lý thuyết mở        & Gemini + local rubric \\
    CODING      & Lập trình           & Gemini + local rubric \\
    \bottomrule
  \end{tabular}
\end{table}
